\documentclass[10pt]{article}

\usepackage[margin=0.82in]{geometry}
\usepackage{amsmath,amssymb}
\usepackage{booktabs,longtable,tabularx,array}
\usepackage[round]{natbib}
\usepackage{microtype}
\usepackage[hidelinks]{hyperref}
\usepackage{xcolor}

\newcommand{\one}{\mathbf{1}}
\newcommand{\good}{\textcolor{green!45!black}{retain}}
\newcommand{\bad}{\textcolor{red!65!black}{reject}}
\newcommand{\diag}{\textcolor{blue!65!black}{diagnostic}}
\newcolumntype{Y}{>{\raggedright\arraybackslash}X}

\title{Bequests and late-life housing exit:\
options, overnight evidence, and the next calibration strategy}
\author{Fertility--Housing project}
\date{15 July 2026}

\begin{document}
\maketitle

\section*{Executive conclusion}

The overnight battery answers a narrower question than ``do households have
bequest motives?'' It asks whether a parent-gated bequest motive becomes useful
after adding one particular late-life housing mechanism: a linear decline in
the owner's borrowing limit from age 66 onward. The answer is no. The
borrowing rule improves the scalar loss only by producing a discrete ownership
collapse between ages 74 and 78, and the valid nested comparison puts the
bequest strength back at zero. This rejects the borrowing rule and the current
combination; it does not reject bequests in a model with a credible old-age
block.

Three additional findings determine the next step. First, the age-82 model
state forces terminal liquidation, so it cannot be compared with the ACS
82--84 ownership bin. Removing that endpoint does not rescue the borrowing
rule: its 74--78 cliff remains. Second, the 12-parameter headline search failed
to recover its own 11-parameter submodel, exposing a search-design failure.
Third, the fresh Jacobian has rank 9 at a relative singular-value tolerance of
$10^{-2}$ and rank 12 only at $10^{-3}$, so the headline system is locally
weakly identified even though the $\theta_0$ column is numerically above the
repeat-noise floor.

The recommended sequence is therefore:
\begin{enumerate}
  \item keep the no-bequest clean frontier as the interim benchmark;
  \item extend the old-age horizon and add externally measured survival before
  any new bequest calibration;
  \item construct HRS/PSID profiles for ownership transitions, downsizing,
  mortgage balances, and nonhousing-wealth decumulation;
  \item use a smooth, data-pinned owner-exit hazard as a low-cost diagnostic,
  then decide from those facts whether the structural model needs medical
  expense risk, a moving/downsizing margin, or a mortgage-balance state;
  \item only after an old-age mechanism passes the ownership and wealth gates,
  reintroduce one free bequest-strength parameter with the restricted winner
  injected into every expanded search.
\end{enumerate}

\section{The decisions and the available options}

There are four distinct decisions. Combining them into one calibration obscures
which margin failed: the form of estate utility, the source of late-life
housing exit, the empirical objects used to discipline that exit, and the
optimizer used to compare nested models.

\subsection{Estate-utility options}

\paragraph{Option B0: no operative bequest motive.}
The clean frontier fixes $\theta_0=\theta_n=0$ and deterministic tenure, leaving
11 estimated coordinates. This is the identified and reproducible interim
benchmark. Its limitation is economic rather than numerical: old-age
ownership remains too high and the model cannot study estate motives or
transfers. Status: \good{} as the benchmark, not as the final old-age model.

\paragraph{Option B1: child-blind luxury warm glow.}
The structural savings literature commonly uses a luxury shifter without a
child-count multiplier \citep{denardi2004,lockwood2018,nakajima2017}. In the
normalization used by the code, the payoff is:
\begin{equation}
v(b)=\theta_0\left[
\frac{(b+\theta_1)^{1-\sigma}-\theta_1^{1-\sigma}}{1-\sigma}
\right].
\label{eq:childblind}
\end{equation}
The subtraction makes zero estate deliver zero utility and prevents a utility
level constant from contaminating fertility comparisons. This is the cleanest
literature baseline once the old-age block is credible. The published
\citet{nakajima2017} calibration reports an annual shifter of \$7,619 in 2000
dollars; there is no verified one-to-one conversion from that number to this
model's income normalization, so $\theta_1$ must remain an explicitly external
scale until a model-specific estate moment is added.

\paragraph{Option B2: parent-gated luxury warm glow.}
The implemented headline variant gives the normalized payoff only to parents:
\begin{equation}
v(b,n)=\one\{n\geq1\}\,\theta_0\left[
\frac{(b+\theta_1)^{1-\sigma}-\theta_1^{1-\sigma}}{1-\sigma}
\right].
\label{eq:parentgate}
\end{equation}
This captures the extensive parent-childless contrast used by \citet{hurd1989}
and the modest difference in motive incidence documented by
\citet{kopczuk2007}. It is less standard than Option B1 because childless
households also leave intended estates in the data. Status: \diag{} as an
extensive-margin sensitivity, not yet the preferred baseline.

\paragraph{Option B3: equal-division robustness.}
If an estate is divided equally across $n$ children, a disciplined alternative
is:
\begin{equation}
v(b,n)=\one\{n\geq1\}\,\theta_0 n\left[
\frac{(b/n+\theta_1)^{1-\sigma}-\theta_1^{1-\sigma}}{1-\sigma}
\right].
\label{eq:equaldivision}
\end{equation}
The form follows the equal-division evidence in \citet{mcgarry1999} and
\citet{wilhelm1996}; it is not a claim that the motive grows linearly with the
number of children. Overnight it returned $\theta_0=0.0106$ and failed the same
ownership test as the headline. Status: \diag{} only.

\paragraph{Option B4: freely estimate both $(\theta_0,\theta_1)$.}
This is attractive economically but not identified by the current target
system. The overnight A5 arm has 13 free parameters, reached
$(\theta_0,\theta_1)=(0.0060,0.1270)$, and failed the ownership rule. Its low
loss is not an estimate because the Jacobian and earlier audits show weak
directions. Status: \bad{} as a calibration, retain only for ridge mapping.

\subsection{Old-age housing and wealth options}

\paragraph{Option E0: the tested borrowing-limit taper.}
The overnight rule multiplies the owner's collateral share by a deterministic
age schedule. It is not a mortgage contract: there is no loan-balance state,
origination date, term, refinancing choice, or scheduled payment. Because the
constraint tightens sharply while the household is still alive, the rule
forces households out of ownership rather than generating gradual repayment.
Status: \bad{}.

\paragraph{Option E1: externally measured survival over a longer horizon.}
Certain death and forced liquidation at age 82 create an artificial endpoint.
Extending the horizon to at least the early 90s and loading age-specific
survival probabilities is the smallest structural correction. It creates
accidental estates and makes the timing of death uncertain, as in
\citet{hurd1989}. Survival alone may not generate enough downsizing, but it is
a prerequisite for interpreting any bequest motive. Status: recommended first
structural change; no calibrated parameter is required.

\paragraph{Option E2: smooth owner-exit or downsizing hazard.}
An age- and state-dependent owner-to-renter or smaller-home hazard can be pinned
from HRS transitions and used as a diagnostic envelope. It would answer whether
a smooth exit path is sufficient before a costly state-space expansion. It is
reduced form and therefore not a final household problem, but it is far safer
than selecting an LTV endpoint by loss. Status: recommended diagnostic after
the data construction.

\paragraph{Option E3: medical expense and health risk.}
Late-life expense risk can generate nonhousing-wealth decumulation and housing
exit while preserving endogenous choice. It is standard in structural
retirement models and interacts meaningfully with intended versus accidental
estates. It requires health states, transition probabilities, and expense
distributions, and it may still leave housing too persistent if moving is
costly. Status: leading structural candidate if the data show health shocks
predict exit.

\paragraph{Option E4: a true mortgage-balance state.}
If the HRS shows that mortgage repayment, refinancing, or equity extraction is
the relevant margin, the model needs a loan balance and contract law of motion.
This is the most faithful household-finance option and the most expensive: it
adds a state and requires origination, term, interest, payment, and refinancing
rules. Reverse-mortgage access can be a later robustness, following
\citet{nakajima2017}, but its low incidence makes it a poor default explanation.
Status: conditional on the data, not the next automatic build.

\paragraph{Option E5: inter vivos down-payment transfers.}
The strongest child-directed evidence concerns transfers while parents are
alive \citep{kvaerner2023,mcgarry1999}. A transfer at family formation would
relax the model's down-payment constraint directly and may be more relevant to
fertility than the terminal estate. This is a separate young-household channel;
it does not repair the late-life ownership block. Status: promising subsequent
extension, identified with gift incidence and amounts rather than estate
moments.

\section{What the overnight battery established}

\subsection{Design}

Every arm re-estimated the 11 clean-frontier coordinates against all 15 active
moments. A1, A3, and A4 added free $\theta_0$; A5 also freed $\theta_1$ and was
labeled unidentified. The owner borrowing endpoint and the headline
$\theta_1$ scale were external variants, never selected by loss. The 22 main
chains completed 14,711 search evaluations, and the 10 fixed-$\theta_0$
profile chains completed 6,143. Every reported point was solved twice at the
tight evaluator with zero repeat difference. Six main-chain and three
profile-chain search winners failed the final market-residual gate and were
excluded.

\subsection{Arm comparison}

Lower loss indicates a closer match to the weighted moments, but it is not an
acceptance criterion by itself. Table~\ref{tab:arms} reports tight results; the
ownership gate is recomputed over ages 62--78 because age 82 is a forced
terminal-liquidation state.

\begin{table}[ht]
\centering
\small
\begin{tabular}{lrrrrcc}
\toprule
arm & free & loss & $\theta_0$ & $\bar h_0$ & 62--78 band & 62--78 cliff \\
\midrule
A0: control & 11 & 6.7881 & 0 & 0.3963 & fail & no \\
A1: parent gate only & 12 & 6.7430 & 0.0211 & 0.3638 & fail & no \\
A2: exit only, $\bar L=0.2$ & 11 & 4.2024 & 0 & 0.2851 & fail & yes \\
A2: exit only, $\bar L=0.4$ & 11 & 5.0115 & 0 & 0.3683 & fail & yes \\
A2: exit only, $\bar L=0.6$ & 11 & 6.1230 & 0 & 0.3818 & fail & yes \\
A3: exit + gate, $\bar L=0.4$ & 12 & 5.2935 & 0.0151 & 0.4237 & fail & yes \\
A3: $\theta_1=0.50$ sensitivity & 12 & 5.6877 & 0.0003 & 0.3783 & fail & yes \\
A4: equal division & 12 & 5.6473 & 0.0106 & 0.3512 & fail & yes \\
A5: both bequest terms free & 13 & 4.9850 & 0.0060 & 0.4148 & fail & yes \\
\bottomrule
\end{tabular}
\caption{Tight arm results. $\bar L$ denotes the externally imposed terminal
multiplier on the owner borrowing limit, not an observed loan-to-value ratio.
A3 at $\bar L=0.2$ produced no tight-eligible chain winner.}
\label{tab:arms}
\end{table}

The scalar gain comes from the wrong lifecycle path. Under primary A2,
ownership at ages 70, 74, and 78 is $0.902$, $0.815$, and $0.043$; under
primary A3 it is $0.861$, $0.755$, and $0.042$. The empirical ownership rate
is roughly three-quarters throughout these ages. The schedule therefore does
not approximate gradual repayment or smooth downsizing.

\subsection{The falsifiable bequest prediction}

A3 contains A2 exactly when $\theta_0=0$. Consequently, a correctly optimized
A3 must weakly improve on the A2 loss of 5.0115. The free A3 chain instead
reported 5.2935, proving that it missed the known nested basin. The dedicated
profile made the same mistake at zero, so the correct profile replaces that
cell with the already verified A2 envelope:

\begin{table}[ht]
\centering
\small
\begin{tabular}{lrrrrr}
\toprule
fixed $\theta_0$ & 0 & 0.05 & 0.10 & 0.20 & 0.40 \\
\midrule
direct profile loss & 8.8134 & 5.0155 & 6.1220 & 8.8669 & 10.7405 \\
valid envelope loss & 5.0115 & 5.0155 & 6.1220 & 8.8669 & 10.7405 \\
\bottomrule
\end{tabular}
\caption{Nuisance-reoptimized $\theta_0$ profile. The zero envelope is the
algebraically equivalent A2 calibration and must be honored in any nested
comparison.}
\label{tab:profile}
\end{table}

Zero weakly dominates the positive grid. The difference between zero and
$0.05$ is only 0.004 loss points, smaller than the uncertainty implied by the
incomplete basin search. The defensible conclusion is therefore ``no evidence
for positive $\theta_0$ under this exit rule,'' not a precise estimate of zero
and not a general rejection of bequests.

\subsection{Identification}

The fresh 15-by-12 Jacobian is deterministic, and the $\theta_0$ column has
target-scaled norm 15.70 against a zero repeat-noise floor. Nevertheless, the
weighted singular values imply rank 9 at relative tolerance $10^{-2}$ and rank
12 only at $10^{-3}$. Thus three parameter combinations are weak at the more
conservative threshold. Since the Jacobian is evaluated at a rejected and
non-optimal A3 point, it is a warning rather than an inference result. Freeing
$\theta_1$ would worsen the problem; A5 remains unidentified by construction.

\section{Plan and decision gates}

The next work should be sequential. Each stage has an observable deliverable
and a stopping rule, so another large calibration cannot conceal a failed
mechanism or optimizer.

\begin{longtable}{@{}p{0.08\textwidth}p{0.25\textwidth}p{0.38\textwidth}p{0.21\textwidth}@{}}
\toprule
stage & action & deliverable & gate \\
\midrule
\endfirsthead
\toprule
stage & action & deliverable & gate \\
\midrule
\endhead
0 & Correct the experiment infrastructure. & Exclude the forced terminal age
from empirical acceptance unless the horizon is extended; require every
expanded optimizer to evaluate and retain the restricted-model incumbent;
seed the zero-profile cell from the A2 winner. & A nested model may never
report a worse best loss than its restricted submodel. \\
\addlinespace
1 & Build the old-age empirical packet before changing the model. & HRS/PSID
profiles with provenance and uncertainty for owner-to-renter transitions,
downsizing, mortgage balances or LTV conditional on ownership, nonhousing
wealth, and the 75+/62--74 decumulation ratio. & No new hard target or
borrowing endpoint without a measured counterpart. \\
\addlinespace
2 & Extend the horizon and add age-specific survival. & A no-bequest solution
through at least the early 90s with the standard diagnostic plots and no
forced endpoint inside the empirical window. & Ownership and wealth paths
must remain numerically stable at the oldest states. \\
\addlinespace
3 & Run a small mechanism frontier with $\theta_0=0$. & Compare a smooth
data-pinned exit hazard with survival plus health/expense risk. Build a true
mortgage state only if the empirical packet points to mortgage balances or
refinancing as the relevant source. & Every four-year nonterminal ownership
bin lies in the predeclared band; no adjacent decline exceeds 15 points; the
wealth path moves in the empirical direction. \\
\addlinespace
4 & Calibrate bequest strength only inside a mechanism that passes Stage 3. &
Estimate the 11 core coordinates plus $\theta_0$, keep $\theta_1$ external,
inject the Stage-3 winner at $\theta_0=0$, and compute a nuisance-reoptimized
profile and fresh Jacobian. & The expanded search recovers the zero incumbent;
positive $\theta_0$ must improve the tight loss materially, pass the lifecycle
gates, and be identified at the declared rank threshold. Otherwise fix
$\theta_0=0$. \\
\addlinespace
5 & Revisit the child channel. & Compare child-blind and parent-gated luxury
warm glow; keep equal division as robustness. Add inter vivos down-payment
transfers only with gift-incidence and amount moments. & Do not infer a child
channel from terminal estates when the identifying variation belongs to
transfers. \\
\bottomrule
\end{longtable}

\section{Recommendation requiring author approval}

The immediate recommendation is to approve Stages 0--2 and defer a new global
calibration. Concretely, keep $\theta_0=\theta_n=0$ in the interim benchmark,
retire the linear borrowing-limit taper, extend the old-age horizon, and build
the HRS/PSID transition and balance profiles. Then run the smooth exit-hazard
diagnostic and survival/expense alternative at $\theta_0=0$. A full
mortgage-balance state should be authorized only if the empirical packet says
the mortgage channel is quantitatively central.

Once one old-age mechanism passes its own lifecycle gates, the preferred
bequest sequence is Option B1 as the literature baseline, Option B2 as the
parent-childless sensitivity, and Option B3 as equal-division robustness. Keep
$\theta_1$ externally fixed until an additional estate-distribution or
bequest-incidence moment identifies it. Inter vivos transfers are a separate,
promising fertility mechanism, but they should follow rather than substitute
for repair of the old-age block.

\clearpage
\appendix
\section{Complete primary target fits}

Tables~\ref{tab:a2fit} and \ref{tab:a3fit} report every active target, model
moment, gap, weight, and loss contribution. They are tight-evaluator records;
the contribution column sums to the reported scalar loss.

\begin{longtable}{p{0.42\textwidth}rrrrr}
\caption{Primary A2: exit rule only, $\bar L=0.4$, tight loss 5.0115.}
\label{tab:a2fit}\\
\toprule
moment & target & model & gap & weight & contribution \\
\midrule
\endfirsthead
\toprule
moment & target & model & gap & weight & contribution \\
\midrule
\endhead
TFR & 1.9180 & 1.9804 & +0.0624 & 20.0 & 0.0778 \\
Childless share & 0.1880 & 0.1966 & +0.0086 & 20.0 & 0.0015 \\
Ownership, ages 30--55 & 0.5755 & 0.6443 & +0.0688 & 100.0 & 0.4733 \\
Parent ownership gap & 0.1677 & 0.1774 & +0.0097 & 45.0 & 0.0043 \\
Housing increment, 0 to 1 child & 0.6644 & 0.6424 & -0.0220 & 14.0 & 0.0068 \\
Old parent-childless nonhousing-wealth gap & 1.0074 & 0.8928 & -0.1147 & 2.0 & 0.0263 \\
Childless renter mean rooms & 3.8053 & 3.8466 & +0.0413 & 6.0 & 0.0102 \\
Childless owner share with at least 6 rooms & 0.5961 & 0.7645 & +0.1684 & 25.0 & 0.7090 \\
Old nonhousing-wealth median & 2.2305 & 1.6781 & -0.5523 & 0.8 & 0.2440 \\
Young childless renter liquid wealth/income & 0.1792 & 0.1808 & +0.0016 & 12.0 & 0.0000 \\
Childless owner-renter mean-room gap & 2.4188 & 2.2787 & -0.1401 & 12.0 & 0.2355 \\
Old-age ownership & 0.7643 & 0.8750 & +0.1107 & 160.0 & 1.9624 \\
Ownership, ages 25--34 & 0.3412 & 0.2630 & -0.0782 & 80.0 & 0.4891 \\
Large-family room gap & 0.3677 & 0.4740 & +0.1063 & 8.0 & 0.0904 \\
Mean occupied rooms & 5.7800 & 5.4431 & -0.3369 & 6.0 & 0.6809 \\
\bottomrule
\end{longtable}

\begin{longtable}{p{0.42\textwidth}rrrrr}
\caption{Primary A3: exit rule plus parent-gated bequests, $\bar L=0.4$,
tight loss 5.2935.}
\label{tab:a3fit}\\
\toprule
moment & target & model & gap & weight & contribution \\
\midrule
\endfirsthead
\toprule
moment & target & model & gap & weight & contribution \\
\midrule
\endhead
TFR & 1.9180 & 2.0048 & +0.0868 & 20.0 & 0.1505 \\
Childless share & 0.1880 & 0.1845 & -0.0035 & 20.0 & 0.0002 \\
Ownership, ages 30--55 & 0.5755 & 0.6047 & +0.0292 & 100.0 & 0.0854 \\
Parent ownership gap & 0.1677 & 0.1424 & -0.0252 & 45.0 & 0.0287 \\
Housing increment, 0 to 1 child & 0.6644 & 0.5323 & -0.1322 & 14.0 & 0.2445 \\
Old parent-childless nonhousing-wealth gap & 1.0074 & 0.9091 & -0.0984 & 2.0 & 0.0194 \\
Childless renter mean rooms & 3.8053 & 4.0932 & +0.2879 & 6.0 & 0.4974 \\
Childless owner share with at least 6 rooms & 0.5961 & 0.8208 & +0.2247 & 25.0 & 1.2623 \\
Old nonhousing-wealth median & 2.2305 & 2.0029 & -0.2275 & 0.8 & 0.0414 \\
Young childless renter liquid wealth/income & 0.1792 & 0.0888 & -0.0904 & 12.0 & 0.0981 \\
Childless owner-renter mean-room gap & 2.4188 & 2.1723 & -0.2465 & 12.0 & 0.7290 \\
Old-age ownership & 0.7643 & 0.8282 & +0.0640 & 160.0 & 0.6544 \\
Ownership, ages 25--34 & 0.3412 & 0.2362 & -0.1050 & 80.0 & 0.8816 \\
Large-family room gap & 0.3677 & 0.3432 & -0.0245 & 8.0 & 0.0048 \\
Mean occupied rooms & 5.7800 & 5.4648 & -0.3151 & 6.0 & 0.5959 \\
\bottomrule
\end{longtable}

\section{Primary A3 parameter table}

The headline arm estimates 12 coordinates and fixes three by external
restriction. Only $\theta_0$ is flagged near its search boundary.

\begin{longtable}{lrrrl}
\toprule
parameter & estimate & lower & upper & role \\
\midrule
\endfirsthead
\toprule
parameter & estimate & lower & upper & role \\
\midrule
\endhead
$\beta$ (annual) & 0.983627 & 0.800 & 0.9995 & estimated \\
$\alpha$ & 0.607209 & 0.020 & 0.980 & estimated \\
$\bar c_0$ & 1.19964 & 0 & 2 & estimated \\
$\bar c_n$ & 0.376812 & 0 & 3 & estimated \\
$\bar h_0$ & 0.423707 & 0.05 & 5.8 & estimated \\
$\bar h_{\mathrm{jump}}$ & 1.41843 & 0 & 8 & estimated \\
$\bar h_n$ & 0.119675 & 0 & 5 & estimated \\
$\psi_{\mathrm{child}}$ & 0.078773 & -3 & 3 & estimated \\
$\kappa_{\mathrm{fert}}$ & 1.83831 & 0.02 & 50 & estimated \\
$\chi$ & 1.09210 & 0.1 & 5 & estimated \\
$H_0$ & 8.06146 & 0.2 & 80 & estimated \\
$\theta_0$ & 0.015074 & 0 & 1.5 & estimated; near zero \\
$\kappa_{\mathrm{tenure}}$ & 0 & \multicolumn{2}{c}{external} & fixed \\
$\theta_1$ & 0.25 & \multicolumn{2}{c}{external scale} & fixed \\
$\theta_n$ & 0 & \multicolumn{2}{c}{external} & fixed \\
\bottomrule
\end{longtable}

\begingroup
\small
\begin{thebibliography}{99}

\bibitem[De Nardi(2004)]{denardi2004}
De Nardi, M. (2004).
Wealth inequality and intergenerational links.
\emph{Review of Economic Studies} 71(3), 743--768.

\bibitem[Hurd(1989)]{hurd1989}
Hurd, M.~D. (1989).
Mortality risk and bequests.
\emph{Econometrica} 57(4), 779--813.

\bibitem[Kopczuk and Lupton(2007)]{kopczuk2007}
Kopczuk, W. and J.~P. Lupton (2007).
To leave or not to leave: The distribution of bequest motives.
\emph{Review of Economic Studies} 74(1), 207--235.

\bibitem[Kv\ae rner(2023)]{kvaerner2023}
Kv\ae rner, J. (2023).
The intergenerational transmission of wealth: Bequest motives and
child-directed transfers.
\emph{Review of Financial Studies}.

\bibitem[Lockwood(2018)]{lockwood2018}
Lockwood, L.~M. (2018).
Incidental bequests and the choice to self-insure late-life risks.
\emph{American Economic Review} 108(9), 2513--2550.

\bibitem[McGarry(1999)]{mcgarry1999}
McGarry, K. (1999).
Inter vivos transfers and intended bequests.
\emph{Journal of Public Economics} 73(3), 321--351.

\bibitem[Nakajima and Telyukova(2017)]{nakajima2017}
Nakajima, M. and I.~A. Telyukova (2017).
Reverse mortgage loans: A quantitative analysis.
\emph{Journal of Finance} 72(2), 911--950.

\bibitem[Wilhelm(1996)]{wilhelm1996}
Wilhelm, M.~O. (1996).
Bequest behavior and the effect of heirs' earnings: Testing the altruistic
model of bequests.
\emph{American Economic Review} 86(4), 874--892.

\end{thebibliography}
\endgroup

\end{document}
